\chapter{Renders}
\label{anx:renders}
\begin{sidewaysfigure}
    \centering
    \includegraphics[width=\textheight, height=0.92\textwidth, keepaspectratio]{Renders/case1.png}
    \caption[Renderizado de red simplificada]{Renderizado de red con visualización simplificada.}
    \label{fig:render_1}
\end{sidewaysfigure}

\begin{sidewaysfigure}
    \centering
    \includegraphics[width=\textheight, height=0.92\textwidth, keepaspectratio]{Renders/case1_2.png}
    \caption[Renderizado de capas detalladas]{Renderizado de red con visualización detallada de capas tipo convolucionales y maxpooling.}
    \label{fig:render_2}
\end{sidewaysfigure}

\begin{sidewaysfigure}
    \centering
    \includegraphics[width=\textheight, height=0.92\textwidth, keepaspectratio]{Renders/case2.png}
    \caption{Renderizado de geometría de cajas.}
    \label{fig:render_3}
\end{sidewaysfigure}

\begin{sidewaysfigure}
    \centering
    \includegraphics[width=\textheight, height=0.92\textwidth, keepaspectratio]{Renders/case3.png}
    \caption[Renderizado de comparación de volumenes]{Renderizado de comparación volúmenes de distintas capas de distintas inferencias.}
    \label{fig:render_4}
\end{sidewaysfigure}

\begin{sidewaysfigure}
    \centering
    \includegraphics[width=\textheight, height=0.92\textwidth, keepaspectratio]{Renders/conexion.png}
    \caption[Renderizado de conexiones]{Renderizado de conexión entre dos capas (operación de convolución).}
    \label{fig:render_5}
\end{sidewaysfigure}